\section{Resumen del capítulo}

\begin{tcolorbox}[resumen,title=Potencias]
Para \(n\in\mathbb N\) con \(n>0\), una potencia representa un producto de factores iguales:
\[
a^n=\underbrace{a\cdot a\cdot\ldots\cdot a}_{n\text{ factores}}.
\]
El número \(a\) es la base, \(n\) es el exponente y el valor de la expresión es la potencia. Además,
\[
a^1=a,
\qquad
a^0=1\quad(a\ne0).
\]
En este capítulo \(0^0\) queda sin definir.
\end{tcolorbox}

\begin{tcolorbox}[resumen,title=Propiedades con exponentes naturales]
Cuando las expresiones están definidas,
\begin{gather*}
a^m\cdot a^n=a^{m+n},\\
\frac{a^m}{a^n}=a^{m-n}\quad(a\ne0,\ m\ge n),\\
(a^m)^n=a^{mn},\\
(a\cdot b)^n=a^n\cdot b^n.
\end{gather*}
Estas propiedades se comprenden contando, cancelando o reorganizando factores iguales.
\end{tcolorbox}

\begin{tcolorbox}[resumen,title=Exponentes enteros negativos]
Si \(a\ne0\) y \(n>0\),
\[
a^{-n}=\frac1{a^n}.
\]
El signo negativo del exponente indica un inverso multiplicativo, no un valor negativo. La expresión \(0^{-n}\) no está definida porque exigiría dividir por cero.
\end{tcolorbox}

\begin{tcolorbox}[resumen,title=Raíces y números reales]
La radicación invierte la pregunta planteada por una potencia. Si \(n\) es impar,
\[
\sqrt[n]{a}=b\Longleftrightarrow b^n=a.
\]
Si \(n\) es par y \(a\ge0\), la raíz principal es el único \(b\ge0\) tal que \(b^n=a\). Por eso,
\[
\sqrt{a^2}=\pm a.
\]
Algunas raíces, como \(\sqrt2\), son reales pero no racionales:
\[
\mathbb N\subseteq\mathbb Z\subseteq\mathbb Q\subseteq\mathbb R.
\]
\end{tcolorbox}

\begin{tcolorbox}[resumen,title=Raíz principal y ecuaciones]
El símbolo radical representa la raíz principal. Así,
\[
\sqrt{25}=5.
\]
En cambio, la ecuación
\[
x^2=25
\]
tiene dos soluciones:
\[
x=\pm5.
\]
La cantidad de soluciones depende de la paridad del exponente y del signo del otro miembro.
\end{tcolorbox}

\begin{tcolorbox}[resumen,title=Exponentes racionales]
Para \(a>0\), \(p\in\mathbb Z\) y \(q\in\mathbb N\) con \(q>0\),
\[
a^{1/q}=\sqrt[q]{a},
\qquad
a^{p/q}=\left(\sqrt[q]{a}\right)^p.
\]
Si \(p<0\), se toma además el inverso multiplicativo. Para una base negativa, \(p/q\) debe reducirse: la potencia es real cuando el denominador reducido es impar y no es real cuando ese denominador es par.
\end{tcolorbox}

\begin{tcolorbox}[resumen,title=Condiciones de las propiedades]
Para \(a,b>0\) y \(r,s\in\mathbb Q\),
\begin{gather*}
a^r\cdot a^s=a^{r+s},\\
\frac{a^r}{a^s}=a^{r-s},\\
(a^r)^s=a^{rs},\\
(a\cdot b)^r=a^r\cdot b^r,\\
\left(\frac ab\right)^r=\frac{a^r}{b^r}.
\end{gather*}
Las restricciones sobre las bases no son detalles separados: forman parte de estas afirmaciones.
\end{tcolorbox}
